\documentclass[10pt]{article}
\usepackage[T1]{fontenc}
\usepackage[utf8]{inputenc}
\usepackage[french]{babel}
\usepackage{amsmath,amssymb,mathrsfs,mathpazo}
\usepackage{geometry}
\geometry{a4paper,margin=2cm}
\usepackage{graphicx,array,tabularx,tabularray,colortbl,multirow,multicol,booktabs,enumitem,pifont,ulem,fancybox,diagbox,cancel}
\usepackage{pstricks,pstricks-add,pst-plot,pst-tree,pst-node,pst-text,pst-func,pst-math,pst-bspline,pst-3dplot,pst-eucl,pst-solides3d}
\let\pscalculate\undefined
\usepackage{pst-bezier}
\usepackage{tikz,pgfplots,tkz-tab}
\usetikzlibrary{arrows,arrows.meta,patterns,calc,positioning,trees,decorations.markings,shapes.geometric,shapes.misc,shapes.arrows,intersections,decorations.pathmorphing,decorations.pathreplacing,backgrounds,shadows,babel}
\pgfplotsset{compat=1.18}
\usepgfplotslibrary{fillbetween}
\usepackage[np]{numprint}
\usepackage{listings,makecell,pas-tableur,twoopt}
\let\barre\undefined
\usepackage[pstricks]{tablvar}
\definecolor{darkgreen}{RGB}{0,100,0}
\providecommand{\red}{\color{red}}
\providecommand{\blue}{\color{blue}}
\providecommand{\green}{\color{green}}
\providecommand{\black}{\color{black}}
\providecommand{\white}{\color{white}}
\providecommand{\gray}{\color{gray}}
\providecommand{\cyan}{\color{cyan}}
\providecommand{\magenta}{\color{magenta}}
\providecommand{\yellow}{\color{yellow}}
\providecommand{\violet}{\color{violet}}
\providecommand{\orange}{\color{orange}}
\providecommand{\euro}{\texteuro}
\providecommand{\EUR}{\texteuro}
\providecommand{\ieme}{\textsuperscript{e}}
\providecommand{\degres}{\ensuremath{^\circ}}
\providecommand{\degree}{\ensuremath{^\circ}}
\providecommand{\vv}[1]{\overrightarrow{#1}}
\providecommand{\Coord}[1]{#1}
\providecommand{\decosix}{}
\providecommand{\starredbullet}{\ensuremath{\bullet}}
\pagestyle{empty}
\setlength{\parindent}{0pt}
\usepackage{adjustbox}\DeclareUnicodeCharacter{2236}{:}\DeclareUnicodeCharacter{2190}{\ensuremath{\leftarrow}}
\providecommand{\R}{\mathbb{R}}
\providecommand{\N}{\mathbb{N}}
\providecommand{\Z}{\mathbb{Z}}
\providecommand{\Q}{\mathbb{Q}}
\providecommand{\D}{\mathbb{D}}
\providecommand{\C}{\mathbb{C}}
\providecommand{\e}{{\mathrm{e}}}
\providecommand{\vect}[1]{\overrightarrow{#1}}
\providecommand{\vectt}[1]{\overrightarrow{\mathrm{#1}}}
\providecommand{\vvt}[1]{\overrightarrow{\mathrm{#1}}}
\providecommand{\np}[2][]{{#2~#1}}
\providecommand{\vv}[1]{{\overrightarrow{#1}}}
\renewcommand{\d}{{\mathrm{d}}}
\providecommand{\degres}{{^\circ}}
\providecommand{\degre}{{^\circ}}
\providecommand{\blue}{\color{blue}}
\providecommand{\red}{\color{red}}
\providecommand{\enskip}{\quad}
\providecommand{\llbracket}{\lbrack\!\lbrack}
\providecommand{\rrbracket}{\rbrack\!\rbrack}
\providecommand{\euro}{{\text{€}}}
\providecommand{\no}{{\text{n°}}}
\providecommand{\interval}[2]{{[#1;#2]}}
\providecommand{\intervFO}[2]{{[#1;#2[}}
\providecommand{\dotfill}{\ldots}

\providecommand{\ds}{\displaystyle}
\providecommand{\textding}[1]{#1}
\providecommand{\Oij}{\(\mathrm{O};\vec{\imath},\vec{\jmath}\)}
\providecommand{\Oijk}{\(\mathrm{O};\vec{\imath},\vec{\jmath},\vec{k}\)}
\providecommand{\Ouv}{\(\mathrm{O};\vec{u},\vec{v}\)}
\providecommand{\intervFF}[2]{\ensuremath{[#1;#2]}}
\renewcommand{\euro}{\ifmmode\text{\texteuro}\else\texteuro\fi}
\DeclareUnicodeCharacter{00B0}{\ifmmode^{\circ}\else\textdegree\fi}
\DeclareUnicodeCharacter{20AC}{\ifmmode\text{\texteuro}\else\texteuro\fi}
% Énoncés : Éducation nationale. Transcription et corrigés : APMEP, Denis Vergès et contributeurs. Découpage et classement : CoopMaths, Sébastien Lozano, Éric Elter et contributeurs. Adaptation de la présentation : Maths Explorer.
% CC BY-SA : https://coopmaths.fr/www/about/
% https://coopmaths.fr/alea/static/bac/2023/tex/bac_2023_03_sujet1_ameriquenord_2_cor.tex
\begin{document}
\noindent\textit{APMEP / CoopMaths. CC BY-SA. Présentation adaptée par Maths Explorer.}\par\medskip
\medskip

\textbf{Partie A - Étude d'une fonction auxiliaire}

\medskip

\begin{enumerate}
\item 
	\begin{enumerate}
		\item Déterminons la limite de la fonction $g$ en $-\infty$.

$\lim\limits_{x \to -\infty} 2x = -\infty$, donc par composition : $\lim\limits_{x \to -\infty} \e^{2x} = \lim\limits_{y \to -\infty} \e^{y} = 0$.

Par ailleurs, $\lim\limits_{x \to -\infty} -2x - 3 = +\infty$

Par limite de somme, on en déduit finalement : $\lim\limits_{x \to -\infty} g(x) = +\infty$.

		\item Pour la limite de la fonction $g$ en $+\infty$, sous la forme actuelle, on a une forme indéterminée.

$\forall x \in \R^*, \quad g(x) = 2x \times \left(3 \times \dfrac{\e^{2x}}{2x} - 1 - \dfrac{3}{2x}\right)$

On a $\lim\limits_{x \to +\infty} 2x = +\infty$.

Or, par composition, $\lim\limits_{x \to +\infty} \dfrac{\e^{2x}}{2x} = \lim\limits_{y \to +\infty} \dfrac{\e^y}{y} = +\infty$, d'après la propriété des croissances comparées.

Par limite de produit, somme et quotient, on a : $\lim\limits_{x \to +\infty} 3\times \dfrac{\e^{2x}}{2x} - 1 - \dfrac{3}{2x} = +\infty$.

Finalement, par limite du produit : $\lim\limits_{x \to +\infty}g(x) = +\infty$.
	\end{enumerate}
\item 
	\begin{enumerate}
		\item Soit $x$ un nombre réel :

$g'(x) = 3 \times 2 \times \e^{2x} - 2\times 1 - 0 = 6 \mathrm{e}^{2 x}-2$.

On arrive bien à l'expression annoncée.

		\item Résolvons l'inéquation suivante :

$\begin{aligned} g'(x) > 0 &\iff 6\e^{2x}-2 > 0\\
&\iff 6\e^{2x} > 2\\
&\iff \e^{2x} > \dfrac{1}{3} \qquad \text{car }6 > 0\\
&\iff 2x > \ln\left(\dfrac{1}{3}\right) \qquad \text{car ln est strictement croissante sur } \R^{*+}\\
&\iff 2x > -\ln(3)\\
&\iff x > \dfrac{-1}{2}\ln(3)
\\ \end{aligned}$

La fonction $g'$ est donc à valeurs strictement positives sur $\left] \dfrac{-1}{2}\ln(3)~;~+\infty\right[$ et à valeurs strictement négatives  sur $\left] -\infty~;~\dfrac{-1}{2}\ln(3)\right[$ (en résolvant une inéquation similaire).

	\item On en déduit le tableau de variations de la fonction $g$ sur $\mathbb{R}$.

Calculons l'image de $\dfrac{-1}{2}\ln(3)$ :

$ \begin{aligned} g\left(\dfrac{-1}{2}\ln(3)\right) &=3\times \e^{2\frac{-1}{2}\ln(3)}-2\times \dfrac{-1}{2}\ln(3) -3 = 3\e^{-\ln(3)} +\ln(3) - 3 = 3\e^{\ln \frac13} +\ln(3) - 3\\&= 3\times \dfrac{1}{3} + \ln(3) - 3 = \ln(3) - 2\\ \end{aligned}$.

		 
		 \begin{center}
		 	\begin{tikzpicture}[scale=0.875]
		 		
		 		\tikzstyle{cadre}=[thin]
		 		\tikzstyle{fleche}=[->,>=latex,thin]
		 		\tikzstyle{nondefini}=[lightgray]
		 		
		 		\def\Lrg{1.5}
		 		\def\HtX{1}
		 		\def\HtY{0.5}
		 		
		 		\def\lignex{-0.5*\HtX}
		 		\def\lignef{-1.5*\HtX}
		 		\def\separateur{-0.5*\Lrg}
		 		
		 		\def\gauche{-1.5*\Lrg}
		 		\def\droite{4.5*\Lrg}
		 		
		 		\def\haut{0.5*\HtX}
		 		\def\bas{-2.5*\HtX-2*\HtY}
		 		
		 		\def\pp#1{\draw(#1,\lignex+0.1*\HtX) -- (#1,\bas);}
		 		\pp{2*\Lrg}

		 		
		 		\node at (-1*\Lrg,0) {$x$};
		 		\node at (0*\Lrg,0) {$-\infty$};
		 		\node at (2*\Lrg,0) {$-\dfrac12 \ln(3)$};
		 		\node at (4*\Lrg,0) {$+\infty$};
		 		
		 		\node at (-1*\Lrg,-1*\HtX) {$g'(x)$};
		 		\node at (0*\Lrg,-1*\HtX) {$$};
		 		\node at (1*\Lrg,-1*\HtX) {$-$};
		 		\node[fill=white] at (2*\Lrg,-1*\HtX) {$0$};
		 		\node at (3*\Lrg,-1*\HtX) {$+$};
		 		\node at (4*\Lrg,-1*\HtX) {$$};
		 		
		 		\node  at (-1*\Lrg,{-2*\HtX+(-1)*\HtY}) {$g$};
		 		\node (f1) at (0*\Lrg,{-2*\HtX+(0)*\HtY}) {$+\infty$};
		 		\node[fill=white] (f2) at (2*\Lrg,{-2*\HtX+(-2)*\HtY}) {$\ln(3)-2$};
		 		\node (f3) at (4*\Lrg,{-2*\HtX+(0)*\HtY}) {$+\infty$};
		 		
		 		\draw[fleche] (f1) -- (f2);
		 		\draw[fleche] (f2) -- (f3);
		 		
		 		\draw[cadre] (\separateur,\haut) -- (\separateur,\bas);
		 		\draw[cadre] (\gauche,\haut) rectangle  (\droite,\bas);
		 		\draw[cadre] (\gauche,\lignex) -- (\droite,\lignex);
		 		\draw[cadre] (\gauche,\lignef) -- (\droite,\lignef);
		 	\end{tikzpicture}
		 \end{center}
		 

		 
		 
		 
		 
		 
		 
		 

La fonction $g$ admet bien un minimum égal à $\ln (3)-2$, atteint pour $x = -\dfrac{1}{2}\ln(3)$.
	\end{enumerate}
\item 
	\begin{enumerate}
		\item Calculons l'image de 0 par $g$ :

$g(0) = 3\e^{2\times 0} - 2\times 0 - 3 = 3 - 0 - 3 = 0$.

On a $g(0) = 0$, donc $x = 0$ est bien solution de l'équation $g(x)=0$.
		\item Le nombre $\ln (3)-2$ est strictement négatif, donc :

\begin{itemize}
\item Sur l'intervalle $\left]-\dfrac{1}{2}\ln(3)~;~+\infty\right[$, où $g$ est strictement croissante, on a $g(0) = 0$, donc $x=0$ est la seule solution à l'équation $g(x) = 0$ sur cet intervalle.
\item Sur l'intervalle $\left]-\infty~;~ -\dfrac{1}{2}\ln(3) \right]$, la fonction $g$ est continue (car dérivable) et strictement décroissante. De plus, 0 est une valeur intermédiaire entre $+\infty = \lim\limits_{x \to -\infty} g(x)$ et $g\left(-\dfrac{1}{2}\ln(3)\right) = \ln(3) - 2 < 0$.

D'après le corollaire au théorème des valeurs intermédiaires appliqué aux fonctions strictement monotones, l'équation $g(x) = 0$ admet une unique solution $\alpha$ dans l'intervalle $\left]-\infty~;~ -\dfrac{1}{2}\ln(3) \right]$.
\end{itemize}

Finalement, l'équation admet donc bien exactement deux solutions sur $\R$ : 

0 et $\alpha$.

Un balayage à la calculatrice donne $-1,5 <\alpha < -1,4$, pour un encadrement au dixième près.
	\end{enumerate}
\item Des variations de $g$ et de la question précédente, on déduit que $g$ est à valeurs strictement positives sur $]-\infty~;~\alpha[$ et sur $]0~;~+\infty[$, et à valeurs strictement négatives sur $]\alpha~;~0[$.
	\end{enumerate}


\textbf{Partie B - Étude de la fonction $f$}

\medskip

\begin{enumerate}
	\item La question posée sous-entend que la fonction $f$ est dérivable sur $\R$. Pour tout $x$ réel, on a :

$f(x) = 3\times \e^{3x} - 2\times \e^x - \left(2x + 1\right)\e^{x} = \e^{x} \times (3\e^{2x} - 2 - 2x-1) = \e^{x}g(x)$.

On arrive bien à l'expression demandée.

	\item La fonction exp est à valeurs strictement positives sur $\R$, donc le signe de $f'(x)$ est le même que le signe de $g$.

Avec la question \textbf{4.} de la \textbf{Partie A}, on peut donc établir que la fonction $f$ est strictement croissante sur $]-\infty~;~\alpha[$ et sur $]0~;~+\infty[$, car sur ces intervalles, sa fonction dérivée est strictement positive, et elle est strictement décroissante sur $[\alpha~;~0]$.

	\item La fonction $f$ n'est pas convexe sur $\mathbb{R}$, car si elle l'était alors sa fonction dérivée $f'$ serait croissante sur $\R$. Or on a vu que la fonction $f'$ est à valeurs positives sur $]-\infty~;~\alpha[$ puis à valeurs strictement négatives sur $]\alpha~;~0[$. On a donc notamment : $f'(-2)>0$ et $f'(-1)< 0$, ce qui prouve que la fonction $f'$ n'est pas croissante sur $\R$ et donc que $f$ n'est pas convexe sur $\R$.
\end{enumerate}

\bigskip

\end{document}
